



\section{Phase 1}\label{sec:phase1}
Here we describe the first phase of the attack, which uses random code shortenings and the GIJS distinguisher to 
recover and unmix all the noisy pairs. Given the $k \times (n+w)$ RLCE public matrix, we proceed as follows.
Let $(n', k', s, d, D, q)$ be GIJS parameters
satisfying the kernel inequality
$(n'-1)\binom{k'+s-1}{s-1} < \binom{k'+d-1}{d}$. Here, $D = k'+w$ is the degree
of the hidden rational curve.

\subsection{Choosing the GIJS parameters}\label{sec:param-choice}

Given the RLCE parameters $(n, k, w, q)$ with public-key length $N = n + w$,
the attacker first fixes a target shortened length $n'$. Shortening at
$N - n'$ positions drops the dimension by the same amount, so
\[
    k' = k - (N - n').
\]
After shortening, the hidden ambient GRS code has dimension $K = k' + w$, so the
hidden rational curve has degree
\[
    D = K - 1 = k' + w - 1.
\]
It remains to choose the jet multiplicity $s$ and the test-polynomial degree
$d$. We search over small integer pairs with $1 \le s \le d < q$ satisfying
both GIJS feasibility inequalities:
\begin{enumerate}
    \item \emph{Reconstruction:} $s(n'-1) > d \cdot D$, which forces every
          $H \in K_\tau$ to vanish at a clean hold-out;
    \item \emph{Kernel:} $(n'-1)\binom{k'+s-1}{s-1} < \binom{k'+d-1}{d}$,
          which guarantees the interpolation kernel is nonzero on random input.
\end{enumerate}
Among feasible pairs $(s,d)$, we prefer smaller $d$: the GIJS matrix has
$\binom{k'+d-1}{d}$ columns, so $d$ dominates the per-run cost. Subject to
that constraint, we prefer a larger relative margin on the kernel inequality,
since a
near-tight kernel is numerically fragile.

There is a tension in choosing $n'$ itself. A smaller $k'$ (hence smaller
$n'$) makes each GIJS run cheaper, but GIJS also needs a hard column-count
floor: $n'$ must be large enough that the reconstruction inequality can hold
at all. In practice $n'$ cannot be pushed far below the boundary where
$s(n'-1) \approx dD$.


\subsection{The Core of the Attack}\label{sec:fullattack}
Let $t'$ denote the number
of noisy columns among the $n'-1$ interpolation columns (all columns except the
holdout column) at any given point in the following algorithm.

\paragraph{Dirty column tolerance.}
We define $t'_{\max}$ to be the maximum number of noisy columns such that PIR
persistence (see Section~\ref{sec:gijs}) guarantees acceptance on a clean
hold-out:
\begin{equation}\label{eq:tmax}
    t'_{\max} = \max\{t' : s\cdot((n'-1) - t') > d\cdot D\}.
\end{equation}
A trial is \emph{recorded} --- its flagged set contributes to the
frequency table --- iff $|V| \leq T_{\mathrm{reject}} := t'_{\max}$.

We deliberately obfuscate an inner step, namely the computation of the set of 
holdout columns $\tau$ on which the GIJS subroutine outputs random. 
While this can be trivially computed by separately running the GIJS subroutine 
as a black box on each holdout column, we have a quicker method, described in 
Section~\ref{sec:reuse-attack}.

\medskip

\begin{algorithm}[H]
    \caption{Phase 1}
    \label{alg:main_analysis}
    \KwIn{Original code parameters $n, w, k, k', q$; raw trial count
        $T$; recording threshold $T_{\mathrm{reject}}$}
    \KwOut{All $w$ pairs $(c_1, c_2)$ of noisy columns in the original code}

    $W \leftarrow \emptyset$\;
    \For{$iter \leftarrow 1$ \KwTo $T$}{
        Choose a random subset $\mathcal{L} \subset [1,n+w]$ where
        $\vert{}\mathcal{L}\vert{} = k - k'$\;
        Compute the $k' \times n'$ shortened code by shortening at
        $\mathcal{L}$\;
        $V \leftarrow \emptyset$\;
        Compute $V$ as the set of holdout columns $\tau \in \{1, \dots n'\}$ on which 
        the GIJS subroutine outputs ``random.'' (see Algorithm~\ref{alg:reuse})\\
        $W \leftarrow W \cup \{V\}$\;
        
    }
    Use $W$ to compute all $w$ pairs of noisy columns (see Algorithm~\ref{alg:cluster})\\
    \Return $w$ twin pairs $(c_1, c_2)$
\end{algorithm}

\subsection{Reuse Speedup}\label{sec:reuse-attack}
We can speed up the attack by noticing that the jet matrix for 
each holdout attack has only small differences from the others. Consider a specific shortening. In the language of the GIJS attack 
we have one shortened point set
$y_1,\dots,y_{n'} \in \Fq^{k'}$. For a hold-out $\tau$, the GIJS
decision is computed from the jet matrix $M_\tau$: its rows are
indexed by $(j,\mathbf{a})$ with $j\neq \tau$ and
$\mathbf{a}\in A_{<s,k'}$, its columns by the monomials
$\mathbf{X}^{\mathbf{m}}$ of degree $d$, and its entry is the Hasse
derivative $(\partial^{[\mathbf{a}]}\mathbf{X}^{\mathbf{m}})(y_j)$.
Writing $z$ for the coefficient vector of $H\in V_d$, we have that $(M_\tau
z)_{(j,\mathbf{a})} = (\partial^{[\mathbf{a}]}H)(y_j)$, so
$\ker M_\tau$ is the interpolation kernel $K_\tau$. Let $\mathbf{v}_\tau
\in \Fq^{N_d}$ be the evaluation vector
$\mathbf{v}_\tau = (y_\tau^{\mathbf{m}})_{\mathbf{m}}$ at which the GIJS 
distinguisher is executed.

Now, let $M_{\mathrm{full}}$ be the jet matrix of \emph{all} $n'$ points,
with $E_F := n'J$ rows grouped in $n'$ consecutive blocks of $J$ rows
(block $j$ holds the jets at $y_j$, with $\mathbf{a}=\mathbf{0}$
first). Note that $M_\tau$ is $M_{\mathrm{full}}$ with block $\tau$
deleted, and that $\mathbf{v}_\tau$ is exactly the
$(\tau,\mathbf{0})$-row of $M_{\mathrm{full}}$, since
$(\partial^{[\mathbf{0}]}\mathbf{X}^{\mathbf{m}})(y_\tau) =
y_\tau^{\mathbf{m}}$.

Call a vector $\mathbf{s} = M_{\mathrm{full}} z$ a \emph{syndrome}. As
$z$ ranges over $\Fq^{N_d}$, the syndromes range over the column
space of $M_{\mathrm{full}}$, which is the kernel of any matrix
$\Pperp$ whose rows form a basis of the \emph{left} kernel of
$M_{\mathrm{full}}$ (the row relations $\mathbf{u}^{T} M_{\mathrm{full}} =
0$). Write $h := E_F - \rk(M_{\mathrm{full}})$ for the number of such
basis rows.

\begin{theorem}[Exactness of hold-out reuse]\label{thm:reuse-exact}
For every $\tau\in[n']$, the following are equivalent:
\begin{enumerate}
\item[(i)] the per-hold-out GIJS decision flags $\tau$;
\item[(ii)] there is a syndrome $\mathbf{s} = M_{\mathrm{full}} z$
supported only on block $\tau$ with $s_{(\tau,\mathbf{0})} \neq 0$;
\item[(iii)] $\rk(\Pperp|_\tau) = \rk(\Pperp|_\tau^{-\mathbf{0}})$,
where $\Pperp|_\tau$ is $\Pperp$ restricted to the $J$ columns of
block $\tau$ and $\Pperp|_\tau^{-\mathbf{0}}$ is the same restriction
with the $\mathbf{a}=\mathbf{0}$ column removed. If $\Pperp$ is empty
($h=0$), then \textup{(iii)} is declared to hold for every $\tau$.
\end{enumerate}
Hence all $n'$ decisions of a trial are obtained from one left-kernel
basis $\Pperp$ and two rank computations per hold-out.
\end{theorem}

\begin{proof}
(i)$\iff$(ii): $\tau$ is flagged iff there is $z\in\ker M_\tau$ with
$\langle z,\mathbf{v}_\tau\rangle \neq 0$. For such a $z$, the
syndrome $\mathbf{s} = M_{\mathrm{full}}z$ vanishes outside block
$\tau$ (that is $M_\tau z = 0$) and
$s_{(\tau,\mathbf{0})} = \langle z, \mathbf{v}_\tau\rangle \neq 0$.
Conversely, if $\mathbf{s} = M_{\mathrm{full}} z$ is supported on
block $\tau$ with $s_{(\tau,\mathbf{0})}\neq 0$, then $M_\tau z = 0$
and $\langle z,\mathbf{v}_\tau\rangle = s_{(\tau,\mathbf{0})} \neq 0$.

(ii)$\iff$(iii): the syndromes supported on block $\tau$ are exactly
the vectors $\mathbf{u}\in\Fq^{J}$ with $\Pperp|_\tau \mathbf{u} = 0$,
i.e.\ $\ker(\Pperp|_\tau)$, embedded in block $\tau$: indeed
$\mathbf{s}\in\ker\Pperp$ iff $\mathbf{s} = M_{\mathrm{full}} z$ for
some $z$. Such a pattern with $u_{\mathbf{0}} \neq 0$ exists iff the
first column of $\Pperp|_\tau$ lies in the span of the remaining
columns: if $\mathbf{u}\in\ker(\Pperp|_\tau)$ has
$u_{\mathbf{0}}\neq 0$, then $\mathrm{col}_{\mathbf{0}} =
-\sum_{i\geq 1} (u_i/u_{\mathbf{0}})\,\mathrm{col}_i$; if
$\mathrm{col}_{\mathbf{0}} = \sum_{i\geq 1} w_i\,\mathrm{col}_i$,
then $(1,-w_1,\dots)\in\ker(\Pperp|_\tau)$ has nonzero first entry.
And $\mathrm{col}_{\mathbf{0}}$ lies in the span of the others iff
removing it does not drop the rank, which is (iii). If $h=0$, then
$M_{\mathrm{full}}$ has full row rank, $\ker\Pperp = \Fq^{E_F}$, the
pattern $\mathbf{e}_{\mathbf{0}}$ is achievable on block $\tau$, and
(ii) holds for every $\tau$ directly.
\end{proof}

Algorithm~\ref{alg:reuse} turns Theorem~\ref{thm:reuse-exact} into a three phase
procedure. In the first
phase it assembles the single matrix $M_{\mathrm{full}}$.
Next, it extracts \emph{every} linear relation
among the rows of $M_{\mathrm{full}}$, in one pass. Forward
elimination is run on the augmented matrix
$[M_{\mathrm{full}} \mid I_{E_F}]$. The identity block
is used as a ledger to record combinations
of the original rows formed by forward elimination. Whenever a row of the
left block has been reduced to zero, the ledger row opposite it is
therefore an explicit combination of the rows of
$M_{\mathrm{full}}$ that cancels them.
Collecting all such ledger rows gives the basis $\Pperp$ of the left
kernel. This elimination step has a cost comparable
to one per-hold-out GIJS run. The third
phase, the loop over $\tau$, merely reads $\Pperp$ locally, one block
at a time.

\begin{algorithm}[H]
\caption{Hold-out reuse: all GIJS decisions of one trial}
\label{alg:reuse}
\KwIn{Shortened columns $y_1,\dots,y_{n'}\in\Fq^{k'}$; GIJS
parameters $(s,d)$.}
\KwOut{The set $V\subseteq[n']$ of flagged hold-outs, identical to
$n'$ separate GIJS runs.}
Build $M_{\mathrm{full}}$ ($E_F \times N_d$), rows grouped in blocks
of $J$, $\mathbf{a}=\mathbf{0}$ first in each block\;
Forward-eliminate the augmented matrix $[M_{\mathrm{full}} \mid
I_{E_F}]$; let $\Pperp$ be the rows of the accumulated transform
opposite the zero rows of the echelon form\;
\If{$h = 0$}{\Return $V = [n']$\;}
$F \leftarrow \emptyset$\;
\For{$\tau \leftarrow 1$ \KwTo $n'$}{
    $A \leftarrow \Pperp$ restricted to the $J$ columns of block
    $\tau$\;
    \If{$\rk(A) = \rk(A \text{ without its first column})$}{
        $V \leftarrow V \cup \{\tau\}$\;
    }
}
\Return $V$\;
\end{algorithm}
\smallskip
Restricting $\Pperp$ to the
$J$ columns of block $\tau$ discards everything a relation says
about the other points and keeps only its entries on $\tau$'s own
rows. According to Theorem~\ref{thm:reuse-exact}, $\tau$ is flagged by the GIJS distinguisher
as random if and only if
the first of these $J$ columns is already a linear
combination of the remaining $J-1$. 
As a result, computing the rank of the matrix of the full $J$ columns and comparing it to 
the rank of the last $J-1$ columns is sufficient; the rank drops if the distinguisher outputs GRS
while it stays the same if it outputs random.

The early exit for
$h = 0$ is a degenerate case: with no
row relations at all, every syndrome is achievable and thus every hold-out
is flagged as random.

\subsection{Unmixing Pairs Via Clustering}\label{sec:cluster}
We now show how to take flagged sets of columns from the tria runs and recover
pairs. Here, the flagged sets of columns mean the set $W$
to which sets of columns flagged as ``random" were added in
Algorithm~\ref{alg:main_analysis}
after each good trial.

The idea is simple once the flagged set of a single trial is looked
at the right way. In a trial shortening, a twin pair can be in only three
states. If both members are kept, then both are dirty columns, so
both are flagged with very very high probability.
If exactly one member is kept, and then the survivor
is derandomized to a clean column per the analysis of the Couvrier et al
paper~\cite{CLT18},
so that column is not flagged, and the
other member is absent from the trial altogether. The final case is where both
members
are hit by the shortening, and neither appears. In all three cases
it is with very very high probability the case that either both are flagged or neither
are flageed.

Finally, a column that belongs to no twin pair, (a pure GRS column), is never
flagged at all.

As a result, if both members of a pair are correctly flagged with high enough
probability,
then the list of trials in
which column A was flagged  is the same as the list of columns for which its
twin column B was flagged.
Meanwhile, it is an empty list for a pure column, and for other twin pairs, the
list is non-empty but with
high probability over the random choices of shortenings distinct. Pair
identification thus becomes grouping
columns by their \emph{pattern}.

\subsubsection{The algorithm}

For each original public column $i$, record its pattern
$b_i \in \{0, 1\}^T$, where $b_i(t) = 1$ if $i$ was flagged in
trial $t$ and $0$ otherwise. Then proceed as follows.

\begin{algorithm}[H]
    \caption{Flag-pattern clustering}
    \label{alg:cluster}
    \KwIn{Flagged sets $V_{1}, \dots, V_{T}$ of $T$ trials,
        mapped back
        to original public-column indices}
    \KwOut{Sets of twin pairs and pure columns; anomalies, if any}
    For each column $i$, form its pattern $b_i$, with $b_i(t) = 1$
    iff $i \in V_t$\;
    Group columns by \emph{exact equality} of patterns\;
    A group of size $2$ with a nonzero pattern is a twin pair\;
    The group with the all-zero pattern is the pure columns\;
    A group of size $1$ with a nonzero pattern is an orphan\;
    A group of size $3$ or more is oversized, and is reported, not
    paired\;
    Compare the size of the all-zero group with the expected
    number $n - w$ of pure columns, and report any excess\;
\end{algorithm}



\subsubsection{Trials for Success}\label{sec:cluster-planning}

We now compute how many trials (assuming that pairs are with overwhelming probability correctly flagged together) 
are necessary to identify and unmix pairs with high probability. Fix a pair. 
The probability both columns of the pair appear in the shortening is \[
    p \;=\; \frac{(n')(n'-1)}{(n+w)(n+w-1)},
\]
the probability that a given pair survives a random shortening to
$n'$ out of $N$ columns. Two distinct pairs separate in a recorded trial
except with probability at
most $p^2 + (1 - p)^2$, when either both pairs appear in the shortening, or neither pair appears in the shortening.


Note that we heuristically treat the pairs as independent here. A union bound over the $\binom{w}{2}$ pairs of
pairs then says that $T$ trials separate every pair from every
other pair with failure probability at most $\epsilon$ as soon as
\[
    \binom{w}{2} \bigl(p^2 + (1 - p)^2\bigr)^T \;\leq\; \epsilon,
\]

and we take $T_{\mathrm{sep}}$ to be the smallest such $T$. A pair
is invisible, and so hides in the pure set, if it is never unhit
in $T$ trials, which happens with probability $(1 - p)^T$ per pair;
a union bound over the $w$ pairs gives the visibility requirement
\[
    w (1 - p)^T \;\leq\; 0.01,
\]
with $T_{\mathrm{vis}}$ the smallest such $T$. Clustering runs $
    T \;=\; \max(T_{\mathrm{sep}}, T_{\mathrm{vis}})
$.

\subsubsection{Empirical Validation}\label{sec:cluster-validation}

Clustering was validated at several points, where the truth
was known from the key and every trial could be checked against
it. 

These parameters include the actual ID-0 parameters. 

\begin{table}[htbp]
    \centering
    \footnotesize
    \begin{tabular}{lcccc}
        \toprule
        Scheme Parameters ($q, n, k, w$) & Attack Parameters ($n',(s,d)$)  & Trials & Flagging exact & Pairs recovered \\
        \midrule
        $(256, 198, 165, 33)$, $(90, 2, 3)$           & $23/23$        & $33/33$    \\
        $(256, 119, 79, 31)$, $(101, 2, 3)$  & $25/25$        & $31/31$   \\
        $(1024, 630, 470, 160)$, $(366, 2, 3)$ & $41/41$ & $160/160$ \\ 
        \bottomrule
    \end{tabular}
    \caption{Cluster validation at reduced parameters. ``Flagging
        exact'' means the flagged set equalled the unhit pairs' columns
        in that many trials.}
    \label{tab:cluster-validation}
\end{table}


\paragraph{ID-0, end to end.} The shipped ID-0 set has since been
run end to end with cluster Phase~1, on a fresh key generated with
the NIST DRBG key generation of the implementation, at the adopted
point $(n', k', (s, d)) = (366, 46, (2, 3))$, with $41$ trials on
$12$ threads. The run succeeded completely: cluster aggregation
over the $41$ trials returned $160$ pairs, $470$ pure columns,
$0$ orphans, and $0$ oversized groups; Phase~2 unmixed $160/160$
pairs, Sidelnikov--Shestakov recovered the correct GRS code, and
verification passed with $160/160$ non-degenerate pairs correct
and $0$ false positives. The run took a little over $3$ hours of
wall time across those $12$ threads to break ID-0. That is the
Phase~1 trial time dominating the run: $41$ trials in batches of
$12$, each trial a single reuse elimination, with aggregation,
Phase~2 (SS bootstrap $0.25$\,s), and verification taking minutes
by comparison. ID-1 and ID-2 have not been run end to end with
cluster Phase~1; their figures below remain model estimates.


